\chapter{Position Independent KV Cache Reuse}
\label{cha:ch2}

\section{Existing Position Independent KV Cache Reuse Schemes}
\label{sec:ch2_intro}

\myx{Same as the last section, you can rename the two subsections as "Existing solutions ..." and "Challenges ..." and exchange their ordering.}

\myx{Also, chapter 3 might benefit from having a title similar to "Latency optimization through cached context"? As we want to transition from the latency challenge to the solution.}

Besides quality, the other major limitation of current code LLM agents is latency (\textbf{Challenge 2}). Many agentic frameworks send a large number of sequential requests to the LLM as it undergoes localization. This iterative process is significantly time-consuming, especially when the prompt in each request includes bulk context information as well as template tokens, and the model re-processes it from scratch every time. For instance, an agent might retrieve a large snippet of code from RepoGraph, feed it into the model to propose a fix, then in a subsequent step feed much of the same snippet again while testing or refining the fix. When no optimizations are prepared for this case, each of those calls recomputes attention over the repeated context tokens, undergoing the same KV cache prefill stage and wasting computation. This redundant computation is a key inefficiency that is shared by many agentic workflows.


Several optimized LLM serving engines have begun to tackle Challenge 2 by introducing caching and scheduling strategies. vLLM \cite{kwon2023vllm} improves throughput by batched execution and paged memory management of the KV cache. It can serve multiple generation requests in parallel and merge requests with identical prefixes so that the shared prefix is computed only once. However, even vLLM only reuses prefix computations when those prefixes exactly align from the start of the prompt. In an agent scenario, calls often share partial contexts that appear after some unique tokens, leading to low cache hit rates. Another approach, Prompt Cache \cite{gim2024promptcachemodularattention}, allows explicitly marking reusable prompt segments and precomputing their KV states. While this yields substantial latency reductions on long-document QA tasks, it requires manual template setup and cannot adapt if templates change.

Cache optimization for LLM inference has recently taken a leap forward with CacheBlend \cite{yao2025cacheblendfastlargelanguage}. CacheBlend’s key innovation is enabling \emph{position-independent} reuse of KV caches. It does so by selectively recomputing a small subset of tokens (HKVD or High-KV-Deviation tokens) to mitigate attention deviation for the new context and updating positional encodings to reflect shifts in token positions. This “blending” enables much higher cache hit rates on retrieval workflows, yielding 2–3× reduction in time-to-first-token and 3–5× throughput improvements, all without degrading output quality. CacheBlend was proved effective in multi-round QA scenarios through datasets such as 2WikiMQA and Musique \cite{ho2020constructingmultihopqadataset, trivedi2022musiquemultihopquestionssinglehop} \mytodo{maybe expand a bit on previous evaluations}. However, there remains a gap in understand how CacheBlend performs in real-world agentic settings, where prompts can consist of long code context and complex repository structure information. In this chapter, we evaluate \myx{maybe "discuss"?} how CacheBlend can reduce latency under agentic use cases. We integrate CacheBlend into a RepoGraph-augmented agent scenario and measure end-to-end latency and throughput against baseline systems like SGLang and vLLM’s (w/ and w/o caching enabled).
